% RESEARCH_DRAFT_SHA256: d4f0c1c5372ef45766fa72b4e6bbebb7e3cfae375c89ba372b588e4ca90f8c7e
% Layout and source-reference wrappers expanded for the web renderer.
\documentclass[11pt]{article}
\usepackage[margin=25mm]{geometry}
\usepackage{fontspec,unicode-math}
\setmainfont{Latin Modern Roman}
\setmathfont{Latin Modern Math}
\usepackage{amsmath,amssymb,mathtools,xurl,hyperref}
\hypersetup{colorlinks=true,urlcolor=blue,linkcolor=blue}
\setcounter{secnumdepth}{0}
\setlength{\parindent}{0pt}
\setlength{\parskip}{5pt}
\setlength{\emergencystretch}{3em}
\newcommand{\sref}[2]{\hyperlink{source:#1:#2}{[S#2]}}
\newcommand{\cataloguescope}[1]{\paragraph{Source record.}#1}
\newcommand{\floor}[1]{\left\lfloor#1\right\rfloor}
\begin{document}
\section*{Research-attempt provenance}
This standalone research attempt retains the complete original definition body below,
including its historical source and preparation notes. Those notes describe the
earlier catalogue preparation; the new attempt and its verification are identified
separately after that body. No original ID, catalogue status, or source statement is changed.

\section{Unique Genome Reconstruction from Seven Beads}\label{statement}

\subsection{Definitions and mathematical statement}

\paragraph{Shared notation.}Standard number systems are $\mathbb N,\mathbb Z,\mathbb Q,\mathbb R,\mathbb C$. Specialized notation and index ranges follow the cited source and the definitions below.


A noiseless reconstruction model assumes its measurements contain no observational error. Unique reconstruction means there is exactly one admissible configuration consistent with those measurements, up to any specified symmetries.
\paragraph{Record scope and qualifications.}
The catalogue title is retained as metadata. Named certificates, contracts, bounds or frameworks in that title are not asserted as results of the original source extract.
The bead constraints, measurement model, and uniqueness symmetries are absent; no positional equations are invented.

Source-record limitation: The bead observation model and the meaning of unambiguous locations are missing; seven known locations alone do not specify a mathematical reconstruction problem.
\paragraph{Statement recorded in the supplied source.}
In the noiseless genome-reconstruction model, do the locations of seven unambiguous beads uniquely determine the locations of all remaining beads? \href{https://doi.org/10.4171/owr/2022/55}{[S2]}
\paragraph{Source status.}
The supplied catalogue labels this source record open; that metadata was checked against the supplied snapshot on 2026-10-11. The label is preserved for provenance and is not an independent verification of present open status or mathematical correctness. The catalogue literature assessment is dated 2026-08-22.
\subsection{Short English statement}
Ask whether seven unambiguously located beads determine the entire stated genome model.
\subsection{Sources}
\begin{itemize}
\item[S1] Ulam AI and the UnsolvedMath Contributors, UnsolvedMath: A Curated Collection of Open Mathematics Problems, catalogue entry. Catalogue curation and metadata: CC BY 4.0. The source extract is retained; definitions, English summaries and qualifications are editorial additions.
\url{https://huggingface.co/datasets/ulamai/UnsolvedMath}.
\item[S2] Algebraic Structures in Statistical Methodology (2023), Oberwolfach Reports, DOI 10.4171/owr/2022/55, pp. 3121--3170
\url{https://doi.org/10.4171/owr/2022/55}.
\item[S3] 3D genome reconstruction primary article
\url{https://pmc.ncbi.nlm.nih.gov/articles/PMC10884149/}.
\item Source-specific Oberwolfach Reports licences, including CC BY-SA 4.0 where applicable, are retained; see THIRD\_PARTY\_NOTICES.md.
\end{itemize}
\paragraph{Preparation and review.}
The original source and UnsolvedMath contributors retain attribution for the source extract. Codex (GPT-6; exact serving revision not exposed) assisted with definitions, English summaries, scope notes and formatting. No new human mathematical authorship or independently verified solution is asserted. Catalogue identity checked against the complete supplied registry and current repository. Distinct upstream records, including repeated questions and broad or damaged extracts, are retained for complete coverage of the source CS category. A repeated record is not a new independent conjecture. Complete mathematical-equivalence and current-literature review has not been performed. Preparation and final review follow contribution prompts C and D, with source gaps disclosed. Structural and rendering checks do not establish a new theorem or open status.
\label{end:statement}

\section{Literature and context check}

I read the primary report on pages 3131--3132 and the later authors' paper,
Cifuentes et al., \emph{3D Genome Reconstruction from Partially Phased Hi-C Data}
(2024), specifically its model equations, Theorem 1 and Conjecture 1.
The report discusses negative rational exponents; the later explicit Conjecture 1
uses exponent $-2$. Theorem 1 supplies generic finite identifiability from six
anchors for negative rational exponents. That theorem is an external dependency,
not re-proved here. The later paper's local ambiguity is exchange of the two
beads in one pair. The argument below respects that ambiguity and does not
invent separately labeled distances from summed contacts.
The seventh-anchor separation route was developed in this attempt. The root
Codex GPT-6 agent suggested a Fubini/semialgebraic bridge to joint genericity;
this agent independently checked it and instead supplied the elementary compact-
projection measurability argument below. This material mathematical contribution
by the root model is disclosed; no human mathematical contribution is asserted.

\section{Research attempt: Finite candidate elimination by a seventh power-law anchor}

\paragraph{A conditional separation theorem for the actual summed observation.}
Fix any real $\alpha<0$ and define the power-law potential of an unordered
pair $P=\{x,y\}$, with multiplicities allowed, by
\[
 F_P(t)=\|t-x\|^\alpha+\|t-y\|^\alpha
 \quad(t\notin\{x,y\}).
\]
Suppose six specified anchors $a_1,\ldots,a_6$ and a true pair $P_*$ have
only finitely many real unordered candidate pairs $P$ satisfying
$F_P(a_i)=F_{P_*}(a_i)$ for all $i$, with every contact finite.
Then for almost every seventh anchor $g\in\mathbb R^3$, the true unordered
pair is the unique candidate satisfying the additional equality
$F_P(g)=F_{P_*}(g)$. For fixed six-anchor data, the good anchors are also an
open dense subset after the finitely many singular locations are removed.
We now prove this, including the nonidentity and exceptional-set claims.

\paragraph{Distinct pairs give different functions of the anchor.}
If two unordered pairs differ, some point $b$ occurs with different
multiplicities, say difference $c\in\{-2,-1,1,2\}$.
As $t\to b$ while avoiding $b$, their potential difference is
\[
 F_P(t)-F_Q(t)=c\|t-b\|^\alpha+O(1).
\]
The remaining terms stay bounded because their centers are different from $b$.
Since $\alpha<0$, the displayed difference diverges in magnitude. Thus it is
not identically zero. On $\mathbb R^3$ minus the finitely many centers it is
real analytic: each squared distance is positive and the real power of a
positive variable has a convergent local power series. That punctured domain
is connected. Indeed two points can be joined by a broken line through an
intermediate point chosen off the finitely many lines through a removed point
and either endpoint; neither segment then meets a removed point.

\paragraph{A nonzero analytic function has a null zero set here.}
For clarity, the needed elementary analytic fact is justified. If an analytic
function on a connected open set vanishes on a nonempty open subset, it vanishes
everywhere: the set where all its derivatives vanish is closed by continuity
and open by its local power series, hence is the whole connected set.
For a nonidentically zero analytic function, cover the domain by countably many
small boxes on which a convergent power series is valid. On each box expand
in its last coordinate about the center. At least one coefficient, an analytic
function of the other coordinates, is not identically zero on the projected box;
otherwise the original function would vanish on that box.
Induct on dimension. In one dimension, a zero is isolated unless every Taylor
coefficient there vanishes; isolated subsets of an interval are countable, by
assigning a rational subinterval isolating each point. In higher dimensions,
the zero set of a selected nonzero coefficient has measure zero by induction.
Outside that exceptional set of the first coordinates, the one-dimensional
slice is not identically zero and has a null zero set. Fubini's theorem gives
measure zero on the box, and the countable cover proves the claim.
The same identity argument shows that each zero set has empty interior; it is
closed relative to the nonsingular domain, so is nowhere dense there.

Apply this fact to $F_P-F_{P_*}$ for each of the finitely many false candidates.
Their zero sets, together with their finitely many centers, form a null set.
A finite union of relatively closed nowhere dense sets is nowhere dense:
inside any open ball, avoid the first by a smaller open ball and repeat for
the finitely many sets. The complement is relatively open and dense.
This proves the conditional separation theorem.

\paragraph{An entirely algebraic special case at $\alpha=-2$.}
Write $q_b(t)=\|t-b\|^2$. For two different pairs $P=\{x,y\}$,
$Q=\{u,v\}$, clearing denominators gives
\[
 p_{P,Q}(t)=(q_x+q_y)q_uq_v-(q_u+q_v)q_xq_y.
\]
This polynomial is not zero: otherwise the two rational functions would agree
away from their poles, contrary to the singularity argument.
Thus for fixed finite candidate data, the bad seventh anchors lie in the zero
set of a finite product of explicit nonzero polynomials, also including the
$q_b$ factors for excluded centers. This proves a nonempty Zariski-open set
of valid seventh anchors for each such fixed data set. It does not on its own
state a uniform polynomial in all the six-anchor data coordinates.

\paragraph{Joint almost-everywhere statement without an unproved projection step.}
Now fix a negative rational $\alpha$ and use the external six-anchor theorem:
outside an exceptional algebraic set of base tuples
$z=(a_1,\ldots,a_6,x_*,y_*)\in\mathbb R^{24}$, the real candidate fiber is
finite. Such a proper algebraic exceptional set is null by the analytic fact
just proved, applied to a nonzero polynomial containing it.
We verify the measurability needed to pass from conditional to joint statements.
Consider tuples $(z,g,x,y)\in\mathbb R^{33}$ satisfying all seven contact
equalities, with $(x,y)$ different from both $(x_*,y_*)$ and $(y_*,x_*)$,
and with all evaluated contacts nonsingular.
For each positive integer $N$, additionally require all coordinates to lie in
$[-N,N]$, each evaluated anchor-to-center distance to be at least $1/N$,
and the distance of the ordered candidate $(x,y)$ from each of the two true
orderings to be at least $1/N$.
The resulting incidence set is compact: these inequalities are closed and
bounded, and the contact functions are continuous with denominators or negative
powers bounded away from their singularities. Its projection onto $(z,g)$ is
compact, by continuity of coordinate projection. Every false nonsingular
candidate appears in one of these compact sets for sufficiently large $N$.
Therefore the set $B$ of bad $(z,g)$ is a countable union of compact sets,
in particular Borel measurable. No general semialgebraic projection theorem
is needed for this step.
For every base tuple in the external full-measure good set, the conditional
separation theorem says its $g$-section of $B$ has measure zero.
Fubini on bounded boxes now gives measure zero for $B$ over those bases.
The excluded null set of bases times $\mathbb R^3$ is also null, by the same
bounded-box argument. Hence for Lebesgue-almost every full seven-anchor and
true-pair tuple, the summed-contact equations uniquely determine the unordered
pair. For a finite number of ambiguous pairs, the conclusion holds simultaneously
outside the finite union of their exceptional sets, when their corresponding
base tuples satisfy the stated generic conditions.

\paragraph{A necessary boundary check.}
Genericity is indispensable. If all anchors lie in the plane $z=0$, reflecting
both beads of a candidate pair across that plane preserves every anchor distance
and hence every summed contact for any negative exponent. Unless the pair
is reflection-invariant up to exchange, the two unordered pairs differ.
For example $\{(0,0,1),(1,0,2)\}$ and
$\{(0,0,-1),(1,0,-2)\}$ give identical data at any number of planar anchors.
This example does not refute the generic statement; it rules out interpreting
``seven'' as a guarantee for arbitrary anchor positions.

\section{Verification and adversarial audit}

The singularity argument counts multiplicities, so coincident candidate beads
cannot slip through it. Finite candidate fibers are an explicit conditional
hypothesis until the external theorem is invoked. The joint argument proves
Borel measurability by compact projections before applying Fubini, rather than
assuming every existential projection is automatically Borel. The generic
conclusion is measure-theoretic for all negative rational exponents; only the
conditional $-2$ result is explicitly expressed by a Zariski-open anchor set.
Swapping the two members of any pair preserves the observations and is never
claimed distinguishable. The proof concerns exact contacts and fixed known
conversion exponent and scale; it supplies no noisy reconstruction algorithm.

\paragraph{Audit corrections.} Recovered the actual summed power-law model and its six-anchor external
theorem instead of substituting ordinary trilateration. Distinguished the report's
negative rational exponents from the later explicit $-2$ conjecture. Supplied the
compact-projection/Fubini bridge for joint almost-everywhere genericity and kept
it distinct from a full-tuple Zariski-open assertion or labeled reconstruction.
\paragraph{External computation.} None; symbolic arguments and boundary checks only.
\section{Scope and remaining gaps}

This establishes a conditional finite-candidate elimination theorem and, using
the explicitly cited six-anchor theorem, an almost-everywhere seven-anchor
identifiability result for unordered pairs for each negative rational exponent.
The abbreviated original record omits the exact genericity convention and the
allowed pair-label and global rigid-motion symmetries. The later paper's
Conjecture 1 asks for sufficiently general full tuples at exponent $-2$;
a uniform Zariski-open full-tuple conclusion is not proved here.
No assertion is made that maternal/paternal labels for every pair can be uniquely
assigned from symmetric sums, nor that noisy data admit a unique exact solution.
These distinctions prevent declaring the entire broad catalogue target solved.
The complete original text and its missing-model warnings remain preserved.

\paragraph{Integration review.} Codex (GPT-6; exact serving revision not exposed) reviewed the written arguments and expanded layout and source-reference wrappers. Repository metadata and the production web renderer were checked using Python 3.13 and JavaScript with MathJax 3.2.2. This model-assisted review is not independent human proof verification; no PDF compilation or formal proof checking is asserted.

\section{Completion Estimate}
\noindent\textbf{COMPLETION: 0\%}
The calculations below the preserved source statement investigate a concrete route,
but establish no substantiated new progress toward the entire original target.
Known special cases and failed sufficient conditions are not credited as a solution.
\section{References}
The original source references and catalogue attribution are preserved in the
Sources section above. The following URLs were actually consulted for this attempt;
retention of an additional original reference does not claim fresh consultation of it.
\begin{itemize}
\item \url{https://ems.press/content/serial-article-files/46992} (consulted 2026-10-11). 
\item \url{https://pmc.ncbi.nlm.nih.gov/articles/PMC10884149/} (consulted 2026-10-11). 
\end{itemize}
\section{Model and contributor disclosure}
Codex GPT-6 completed written-proof review of this attempt at its declared scope; this is model-assisted review, not independent human proof checking.
OpenAI gpt-6-astra, requested reasoning configuration ultra, generated this research
attempt on 2026-10-11. The exact serving revision is not exposed. Codex GPT-6
performed orchestration and repository-format review and suggested a Fubini/projection bridge for joint genericity; its exact serving revision
is also not exposed. The mathematical route, boundary checks and final adversarial
self-audit followed contribution prompts A then B in
\texttt{docs/contribution\_prompts.md}, after reading \texttt{CONTRIBUTING.md}.
Python 3.14.2 assembled the standalone documents; any mathematical computation is
specified in Verification. No theorem prover was used. No human mathematical
authorship, substantial human proof correction, or independent human proof checking
is asserted. Model-assisted review, source retrieval, and structural checks do not
establish correctness or novelty. Full private reasoning transcripts are unavailable
as reproducibility artifacts; the complete arguments, inputs and outcomes actually
used are stated here or in the companion batch audit.
\section{Final status}
Unresolved as a new research contribution. This is a concrete route investigation,
not a claim of new complete resolution or a current-status certification. A known
answer reproduced for context does not become a new contribution. The catalogue's
source-status label is preserved; this attempt status does not assert that a
source-labelled solved question is open.
\end{document}
